\section{Algorithms for Incremental Discovery of Event Queries}
\label{sec:algos}

In this section, we show how event query discovery is
realized
incrementally, under the changes discussed in the previous section. To this
end, we first summarize a basic discovery procedure that serves as a
baseline (\autoref{sec:discovery}). Then, we discuss how
to update
the set of matching queries (i) once streams are added
and/or the support threshold is increased (\autoref{sec:insert_increment}),
and (ii) once
existing streams are deleted and/or the support threshold
is decreased (\autoref{sec:delete_decrease}).

\subsection{Basic Algorithm for Event Query Discovery}
\label{sec:discovery}

As a basic algorithm for event query discovery, we consider an approach that
searches through the space of candidate queries to collect all matching
queries. To this end, it follows an iterative generate-and-test scheme. For
each query that matches the stream database, stricter child queries are
generated through the addition of query terms, or the insertion of variables
and attribute values for placeholders in existing
query terms. The child
queries are, in turn, evaluated over the stream
database. If they match, they become part of the
result set and their child
queries are considered further.

By starting the approach with an empty query, the approach iteratively
explores the search space of candidate queries. It terminates once all child
queries of all matching queries have been assessed.

The essence of this approach is captured by the function \textsc{Discovery},
formalized in \autoref{alg-discovery}.
It takes as input a stream database $D$; a support threshold $\theta$;
the frequent stream alphabet $\Gamma(D,\theta)$;
a set of queries $Q$ that are frequent in $D$; a parent dictionary
$P$ that is modelled as a
function $P: \mathcal{Q} \to \mathcal{Q}$, mapping child queries to their
parent; and a stack $O$ of queries that are frequent in $D$ and still shall
be explored by a depth-first search through the
space of candidate queries.

\begin{algorithm}[t]
	\small
	\caption{\textsc{Discovery}}
	\label{alg-discovery}
	\KwIn{\, \ Stream database $D$,
		%event schema $\mathcal{A}$, sequence of query var. $\mathcal{X}^<$,
		support threshold $\theta$,
		frequent stream alphabet $\Gamma(D,\theta)$,
		 \\ \hspace{3em} \,
		set of queries $Q$, parent dictionary $P$, query stack $O$
		% \\ \hspace{3em} \,
		% Boolean flag $b\in \{\True,\False \}$ whether to insert variables
		.}
	\KwOut{Updated set of queries $Q$, updated parent dicitionary $P$.}
	\BlankLine

	\tcc{\scriptsize Explore queries that are pushed to the stack}
	%${O \leftarrow}$ empty stack; $O.\textsc{push}(q)$\;
	\While{${O}$ not empty}
	{
		${q \leftarrow O}.\textsc{pop}()$\;
		\If{\normalfont{\textsc{Match}}$({q,D,\theta}$)} {
			\tcc{\scriptsize For frequent queries,
				explore their children}
			${Q \leftarrow Q \cup \left\{q\right\} }$\;
			${P \leftarrow}$
			\textsc{ChildQueries}($\mathit{q,\Gamma(D,\theta),P}$)\;
			${C} \leftarrow  \{q'\in \dom(P) \mid P(q')=q\}$\;
			\If{$ C \neq \emptyset$}{
				\lForEach{$q'\in C$}{
					$O.\textsc{push}(q')$}
			}
		}
	}
	{\Return{$\mathit{Q,P}$}}
\end{algorithm}

Function \textsc{Discovery} then explores the queries on the given stack.
Each of them is evaluated against the stream database using the Boolean
function \textsc{Match}, which returns true if query $q$ is frequent in the
stream database $D$ as defined by the support threshold $\theta$; otherwise,
it returns false. If the query is frequent, it is added to the result set
$Q$ and its child queries are generated using function \textsc{ChildQueries}
and recorded in the aforementioned parent dictionary $P$. If such child
queries have indeed been generated, i.e., set $C$ is not empty, these
queries are pushed to the stack for further exploration.
When the stack is empty, the function returns the set of matching
queries as well as the updated parent dictionary.

Moreover, we summarize the functionality of \textsc{ChildQueries} to
generate child queries of a query $q$ as follows. It generates child queries
by the insertion of a \emph{single} symbol, i.e., a variable or an attribute
value from the frequent stream alphabet, for a placeholder in one of the
existing
query terms. In addition, it
also generates further child queries by appending a new query term to the
query that is empty (i.e., comprises placeholders) except for one position
that is filled with a symbol (variable or attribute value). In general,
these insertions consider all possible attribute values, as well as a set of
possible variables (for which the size is bounded by half of the length of
the shortest stream in the database).

While we use function \textsc{Discovery} also in our incremental approaches
for event query discovery, a basic discovery algorithm simply instantiates
by adding the empty query to the query stack and creating its entry in the
parent dictionary. We formalize this instantiation in
\autoref{alg-baseline}. Given a stream database $D$ and support threshold
$\theta$, it initializes the result set $Q$, derives the frequent stream
alphabet $\Gamma(D,\theta)$, constructs the empty query
$\langle\varepsilon\rangle$ and pushes it to the stack, before starting the
exploration with function \textsc{Discovery}.

\begin{algorithm}[H]
	\small
	\caption{\textsc{Basic Discovery Algorithm} }
	\label{alg-baseline}
	\KwIn{\, \ Stream database $D$, support threshold $\theta$.}
	\KwOut{Set of matching queries $Q$, updated parent dictionary $P$.}
	\BlankLine
	$\mathit{Q \leftarrow \emptyset}$\;
	$\Gamma(D,\theta) = \{a \in
	\Gamma \mid |\{S \in D \mid \exists \ j\in\{1,\dots,|S|\},
	i\in\{1,\dots,n\}: S[j].A_i = a\}| \geq s(|D|,\theta)\}$\;
%	$\Gamma_D \gets \{a \in \Gamma \mid |M_+(a,D)| \geq \lceil |D| \cdot
%	\theta \rceil\}$\;
	$q \gets\langle\varepsilon\rangle$\;
	$P(q) \gets
	\langle\varepsilon\rangle$\;

	${O \leftarrow}$ empty stack\;
	$O.\textsc{push}(q)$\;
	$ Q,P \leftarrow \textsc{Discovery}(D, \theta, \Gamma(D,\theta), Q, P,
	O)$\;
	{\Return{$\mathit{Q,P}$}}
\end{algorithm}

%The baseline algorithms \autoref{alg-baseline} relies
%on an iterative generation of query candidates,
%which are then matched against a stream database. It follows a
%bottom-up direction in the exploration of candidate queries and step-wise
%generation of stricter queries for a given query $q$. It starts with the
%empty
%query, which will be elaborated further later.
%The child queries of $q$ are those obtained by inserting
%a new variable, by inserting a variable that has been present already, or by
%inserting an attribute value; either way, considering an existing query term
%or a newly added query term.
%\color{blue}
%Variables are inserted in a certain order, which
%is reflected in a total order $<$ over a set of variables $\mathcal{X}$. To
%simplify the notation, we write $\mathcal{X}^<$ for the sequence of
%variables induced by $<$ over $\mathcal{X}$. Also, our
%construction employs at most $s/2$ variables with $s=\min_{S\in D} |S|$
%being the minimal stream length in the stream database. Hence, the size of
%$\mathcal{X}$ and the length of $\mathcal{X}^<$ are bounded.
%
%Query candidate generation, formalized in
%\autoref{alg-bu-df:childquery}, takes as input a query $q$; an event
%schema $\mathcal{A}$; a supported alphabet $\Gamma_D$; a parent dictionary
%$P$ that is modelled as a
%function $P: \mathcal{Q} \to \mathcal{Q}$, mapping child queries to their
%parent; a sequence of variables $\mathcal{X}^<$; and a Boolean flag $b$ to
%control whether variables shall be inserted. The algorithm returns the set
%of child queries for $q$.
%
%\autoref{alg-bu-df:childquery} is divided into three parts, following an
%initialization (\autoref{alg1:init_start} - \ref{alg1:init_end}),
%each capturing one type of insertion. It is designed such that \emph{all}
%child queries are obtained by
%insertion
%of a single symbol (attribute value or variable) are constructed
%\emph{exactly once}. The generation relies on the auxiliary function
%$\textsc{Next}$
%(\autoref{alg1:next_start} - \ref{alg1:next_end}). Given a query, it
%generates query candidates by inserting a given symbol (variable or
%attribute value) in
%any query term following a given position in the query, if the query
%contains a placeholder for that attribute. In addition, it
%generates candidates by appending a new query term to the query that
%is empty except for the given symbol for the given attribute. Based thereon,
%the three parts of the main algorithm include:
%
%\textit{Insertion of a new variable} (\autoref{rule:insertnewvar} -
%\ref{rule:insertnewvar_end}). For each attribute, we consider all query
%terms after the last occurrence of any attribute value ($g$) and after the
%first
%occurrence of the last inserted variable ($f$). Using function
%$\textsc{Next}$, a single new variable $\mathcal{X}^<[v+1]$ is introduced
%by replacing a placeholder or adding a new query term.
%
%\begin{algorithm}[H]
%	\footnotesize
%	\caption{\textsc{ChildQueries}: Query Cand. Generation}
%	\label{alg-bu-df:childquery}
%	\KwIn{\, \ Query $q$, event schema $\mathcal{A}$, supported alphabet
%		$\Gamma_D$, \\ \hspace{3em} \,
%		parent dictionary $P$, sequence of query variables $\mathcal{X}^<$,
%		\\ \hspace{3em} \,
%		Boolean flag $b\in \{\True,\False \}$ whether to insert variables.
%	}
%	%\tcc*{some comment}		\;
%	\KwOut{Updated parent dictionary $P$.}
%	\SetKwFunction{FNext}{\textnormal{\textsc{Next}}}%
%	\BlankLine
%	$\mathit{g \gets}$ position of last inserted attribute value in $q$\;
%	\label{alg1:init_start}
%	$\mathit{f \gets}$ first position of last inserted variable in $q$\;
%	$z \gets \max(g,f)$\;
%	$\mathit{n \leftarrow} |\mathcal{A}|$\;
%	\label{alg1:init_end}
%	\tcc{\scriptsize Shall variables be inserted?}
%	\If{$b= \True$}{
%		$v \gets$ number of distinct query variables in $q$\;
%		\tcc{\scriptsize Insert a new variable
%			$\mathcal{X}^<[v+1]$}
%		\ForEach{$A\in\mathcal{A}$}{
%			\label{rule:insertnewvar}
%			\ForEach{$q'\in
%				\FNext(q,z,A,\mathcal{X}^<[v+1],n)$}{
%				\label{rule:insertnewvar_step2}
%				\ForEach{$q''\in
%					\FNext(q',z+|q'|-|q|,A,\mathcal{X}^<[v+1],n)$}{
%					$P(q'') \leftarrow q$\;
%				}
%			}
%		}
%		\label{rule:insertnewvar_end}
%		\tcc{\scriptsize Insert last inserted variable
%			$\mathcal{X}^<[v]$}
%		$l \gets$ last position of last inserted variable in $q$\;
%		$s \gets$ last inserted symbol in $q$\;
%
%		\If{$s=\mathcal{X}^<[v]$}{
%			\label{rule:insertexvar}
%			$\mathit{A \leftarrow}$ attribute $A\in \mathcal{A}$ with $s\in
%			\dom(A)$\;
%			\lForEach{$q'\in \FNext(q,l,A,\mathcal{X}^<[v],n)$}{
%				$P(q') \leftarrow q$
%			}
%		}
%		\label{rule:insertexvar_end}
%	}
%	\tcc{\scriptsize Insert new supported attribute value}
%	\ForEach{$A\in\mathcal{A}$}{
%		\label{rule:inserttype}
%		\ForEach{$a\in (\Gamma_{D}\cap \dom(A))$}
%		{
%			\lForEach{$q'\in \FNext(q,z,A,a,n)$}{
%				$P(q') \leftarrow q$
%			}
%		}
%	}
%	\label{rule:inserttype_end}
%	\Return{$\mathit{P}$}
%
%
%	\BlankLine
%
%	\tcc{\scriptsize Function to generate the next candidate queries}
%	\FNext{\textnormal{query} $q$, \textnormal{start pos.} $s$,
%		\textnormal{attribute} $A$, \textnormal{symbol}	$y$,
%		\textnormal{number
%			of attributes} $n$}{
%		\label{alg1:next_start}
%
%		$\mathit{Q' \leftarrow \emptyset}$\;
%		\ForEach{$i\in \{s,\ldots, |q|\}$}
%		{
%			\tcc{\scriptsize Generate candidate by replacing placeholder
%				with
%				given variable/attribute value $y$}
%			\If{$q[i].A=\_$}{
%				$q' \leftarrow  q$\;
%				$q'[i].A \leftarrow y$\;
%				$Q' \leftarrow Q' \cup \left\{q'\right\}$\;
%
%			}
%			\tcc{\scriptsize Generate candidate by including new query term
%				that
%				only contains given variable/attribute value $y$}
%			$e' \leftarrow \varepsilon$\;
%			$e'.A \leftarrow y$\;
%			$q'\leftarrow \langle q[1],\dots,
%			q[i], e', q[i+1],\dots, q[|q|]
%			\rangle$\;
%			$Q' \leftarrow Q' \cup \left\{q'\right\}$\;
%		}
%		\Return{$\mathit{Q'}$}
%		\label{alg1:next_end}
%	}
%\end{algorithm}


\subsection{Stream Insertion and/or Increased Support
Threshold}
\label{sec:insert_increment}

Having discussed a basic discovery algorithm, we now
present a first algorithmic solution to incremental
processing under changes. Specifically, we consider the
case that streams are added and/or that the support
threshold is increased.

We summarize the idea behind the algorithm as follows:
From
\autoref{lemma:update}, we know that the set of matching
queries after the update will be a subset of the original
set of matching queries. As such, no new queries need to
be generated; we only need to check whether the
queries in the given set are still
frequent. Moreover, the scope of this verification depends
on the support threshold. If it remains unchanged and only
the database changed, it suffices to evaluate the existing
queries over the newly added streams. Otherwise, if
the threshold is increased as well, the entire updated
stream database needs to be considered.

We formalize this approach in \autoref{alg_gen}. It
takes as input the original stream database $D_0$; the
original support threshold $\theta_0$; the
set of matching queries $Q_0$ obtained for $D_0$ and
$\theta_0$; the parent dictionary $P$ for these
queries; the updated stream database $D_1$; and the
updated support threshold $\theta_1$. It returns the
set of matching queries $Q_1$ for $D_1$ and
$\theta_1$ along with the updated parent dictionary
$P$.

After initializing the result set $Q_1$, the
algorithm determines the scope for the verification
of the existing queries, as the entire updated
stream database $D_1$ or the added streams
$D_1\setminus D_0$. Then, the verification is
operationalized with stack onto
which we push the queries that shall be checked. We start
with the empty query and then go through $Q_0$ by means of
the parent-child-relation between queries, as stored in
the parent dictionary.

%If a queries matches, it is added to the updated query
%set and its children are
%added to the stack of queries to be checked for their
%matching behavior.

% In case it does not match, we we add its parent query to the stack of queries to be checked.
% Then discard the query entry from the parent dictionary $P$.

%Once the stack is empty, the algorithm returns the
%updated query set $Q_1$ and the parent dictionary $P$.


\begin{algorithm}[t]
    \small
    \caption{\textsc{UpdatePlus}}
    \label{alg_gen}
    \KwIn{\, \ Stream database $D_0$, support threshold $\theta_0$,
    set of matching queries $Q_0$,
	 \\ \hspace{3em} \,
    parent dictionary $P$,
    updated stream database $D_1$,
    updated support threshold $\theta_1$.}
    \KwOut{Set of matching queries $Q_1$, parent
    dictionary $P$.}
    \BlankLine
    $\mathit{Q_1 \leftarrow \emptyset}$\;
    \lIf{$\theta_1 > \theta_0$}
    {
        $D \leftarrow D_1$}
    \lElse{$D \leftarrow D_1\backslash D_0$}

%    $\mathit{R \leftarrow \emptyset}$\;
    ${O \leftarrow}$ empty stack\;
    % \lForEach{$q\in Q$}
    %     {$R \leftarrow R \cup \left\{P(q)\right\} $}
    % \lForEach{$q\in Q\backslash R$}{
    % $O.\textsc{push}(q)$}
    $O.\textsc{push}(\langle\varepsilon\rangle)$\;
    \tcc{\scriptsize Explore queries that are pushed to the stack}
    \While{${O}$ not empty}
    {
            ${q \leftarrow O}.\textsc{pop}()$\;
            % \If{$q \not\in Q_1$}{

%            $\mathit{M \leftarrow}$
%\textsc{Match}(${q,D,\theta}$)
%            \;
            % \If{$ M = \False$} {
            \If{\normalfont{\textsc{Match}}(${q,D,\theta}$)}
             {
            \tcc{\scriptsize For frequent queries, add
            the query
            to the result set and
            push their child queries onto the stack}

            ${Q_1 \leftarrow Q_1 \cup \left\{q\right\}
            }$\;
            % $O.\textsc{push}(P(q))$\;
            % $P.\textsc{pop}(q)$ \;
            ${C} \leftarrow  \{q'\in \dom(P) \mid P(q')=q\}$\;
            \lForEach{$q'\in C$}{
            $O.\textsc{push}(q')$}
            }
            % \lElse{
            %     ${Q_1 \leftarrow Q_1 \cup \left\{q\right\} }$}

        % }
    }
    {\Return{$\mathit{Q_1,P}$}}
    \end{algorithm}

\subsection{Stream Deletion and/or Decreased Support Threshold}
\label{sec:delete_decrease}

Next, we turn to incremental event query discovery once
streams are deleted from the database and/or the support
threshold is decreased. Again, our approach relies on
\autoref{lemma:update}, and exploits that all queries in
the existing set of matching queries will
also match the updated stream database.

In general, the idea of handling stream deletion and/or a decreased support
threshold in an incremental manner can be summarized as follows.
We consider the queries, for which the child queries are not frequent.
Intuitively, these queries denote the frontier of frequent queries in the
space of all query candidates. For these queries, we derive the child
queries and add them to the stack of queries to be explored further using
the basic discovery procedure. However, this exploration first considers
solely the attribute values that have been frequent in the original
database. In order to ensure completeness of the set of matching queries, in
a second step, we therefore also check whether the frequent alphabet in the
updated stream database contains new attribute values. If so, we generate
the child queries of all frequent queries found so far, but only incorporate
the new frequent attribute values (and no insertion of variables) in the
respective generation. All queries generated in this way, had not
been considered previously. If one of these queries is frequent, we
iteratively continue the exploration with their child queries, this 
time inserting all frequent attribute values as well as variables.
As such, the algorithm is similar to the basic discovery procedure
introduced in \autoref{sec:discovery}, but shows an important difference: It
only assesses queries that have been not been checked as part of the
previous discovery run over the database before the change.



%For these queries, we construct the child
%calculating their children and adding them to the stack of queries to be
%matched.
%Then the algorithm checks whether the supported alphabet of the new stream
%database is the same
%as the supported alphabet of the old stream database.
%In case, the supported alphabet includes new types, we consider all
%matching queries calculated so far
%and generate new candidate queries by adding the new types to the supported
%alphabet.
%We prevent queries from being generated more than once by first adding to
%the existing queries
%only the newly supported types. All queries generated in this way, had not
%been considered
%previously.
%If a query matches, we keep generating
%(then considering the entire supported alphabet) new candidate queries
%until we reach a non-match.
%% Finally, we disgard all matching queries which are not descriptive to
%%obtain the updated set of descriptive queries.
%The algorithm is similar to the baseline algorithm with the main difference
%that it only checks queries that were not part of the previous event query
%discovery.


We formalize the respective algorithm in \autoref{alg_spec}. The input is
the same as discussed for \autoref{alg_gen}, in terms of the stream
databases, support thresholds, parent dictionary and the set of matching
queries. The algorithm then initializes the frequent alphabets for the
stream databases before and after the change, i.e., $\Gamma(D_0,\theta_0)$
and $\Gamma(D_1,\theta_1)$. Then, we collect all parents of frequent
queries, as a set $R$, to be able to iterate over all queries, for which the
child queries are not frequent, given by $Q_1 \setminus R$. For these
queries, the child queries are generated using function
\textsc{ChildQueries}, considering the frequent alphabet
$\Gamma(D_0,\theta_0)$, and pushed to a query stack. The latter is used in
function \textsc{Discovery}, from \autoref{alg-discovery}, to explore the
space of candidate queries, starting from the queries on the stack.

\begin{algorithm}[t]
    \small
    \caption{\textsc{UpdateMinus}}
    \label{alg_spec}
    \KwIn{\, \ Stream database $D_0$, support threshold $\theta_0$,
    set of matching queries $Q_0$,
    \\ \hspace{3em} \,
    parent dictionary $P$,
    updated stream database $D_1$, updated support
    threshold $\theta_1$.}
    \KwOut{Set of matching queries $Q_1$, updated parent
    dictionary $P$.}
    \BlankLine
    $\mathit{Q_1 \leftarrow Q_0}$\;
    $s_0 \gets s(|D_0|,\theta_0)$\;
    $\Gamma(D_0,\theta_0) \gets \{a \in
    \Gamma \mid |\{S \in D_0 \mid \exists \
    j\in\{1,...,|S|\},
    i\in\{1,...,n\}: S[j].A_i = a\}|\geq s_0\}$\;
    $s_1 \gets s(|D_1|,\theta_1)$\;
    $\Gamma(D_1,\theta_1) \gets \{a \in
    \Gamma \mid |\{S \in D_1 \mid \exists \
    j\in\{1,...,|S|\},
    i\in\{1,...,n\}: S[j].A_i = a\}| \geq s_1\}$\;
%    \Gamma_{D_0} \gets \{a \in \Gamma_0 \mid |M_+(a,D_0)|
%\geq \lceil |D_0| \cdot \theta_0 \rceil\}$\;
%    $\Gamma_{D_1} \gets \{a \in \Gamma_1 \mid
%|M_+(a,D_1)| \geq \lceil |D_1| \cdot \theta_1 \rceil\}$\;
    $\mathit{R \leftarrow \emptyset}$\;
    ${O \leftarrow}$ empty stack\;
    \lForEach{$q\in Q_1$}
        {$R \leftarrow R \cup \left\{P(q)\right\} $}
    \ForEach{$q\in Q_1\backslash R$}{
        ${P \leftarrow}$
        \textsc{ChildQueries}($\mathit{q,\Gamma({D_0},\theta_0),P}$)\;
        ${C} \leftarrow  \{q'\in \dom(P) \mid P(q')=q\}$\;
        \If{$ C \neq \emptyset$}{
        \lForEach{$q'\in C$}{
            $O.\textsc{push}(q')$}
        }
    }
    $ Q_1,P \leftarrow \textsc{Discovery}(D_1,\theta_1,
    \Gamma({D_0},\theta_0), Q_1, P, O)$\;

    \If{$\Gamma({D_1},\theta_1)\backslash
    \Gamma({D_0},\theta_0) \neq \emptyset$}{
        \ForEach{$q\in Q_1$}{
        ${P \leftarrow}$
        \textsc{ChildQueriesNoVars}($\mathit{q,\Gamma({D_1},\theta_1)\backslash
         \Gamma({D_0},\theta_0),P}$)\;
        ${C} \leftarrow  \{q'\in \dom(P) \mid P(q')=q\}$\;
        \If{$ C \neq \emptyset$}{
        \lForEach{$q'\in C$}{
            $O.\textsc{push}(q')$}
        }
    }
    $Q_1,P \leftarrow \textsc{Discovery}(D_1,
    \theta_1, \Gamma({D_1},\theta_1), Q_1, P, O)$\;

    }
    {\Return{$\mathit{Q_1,P}$}}
    \end{algorithm}

As mentioned, we need to consider that an update to the stream
database changed the frequent alphabet, i.e., attribute values became
frequent due to the removal of streams. If so, i.e., if
$\Gamma({D_1},\theta_1)\backslash
\Gamma({D_0},\theta_0)$ is not empty, for each query obtained so far, we
need to consider the child queries constructed
only through the insertion of these new frequent attribute values
(neglecting the insertion of additional variables), which is captured by
function \textsc{ChildQueriesNoVars}. For these queries, further exploration
may be needed, so that they are pushed onto the query stack, which is then
used in function \textsc{Discovery}.
